% topic: algorithms
\question \textbf{การดำเนินการทางตรรกะระดับบิต (Bitwise Operations)} คือการดำเนินการทางตรรกศาสตร์กับเลขฐานสองทีละหลัก (บิต) โดยมีตัวดำเนินการพื้นฐาน 4 ชนิด สำหรับเลขฐานสองขนาด \textbf{4 หลัก} ดังนี้:

\medskip
\noindent
กำหนด $X = 1011_2$ และ $Y = 0110_2$

\begin{itemize}
  \item \textbf{AND (\&):} ผลลัพธ์เป็น 1 ก็ต่อเมื่อทั้งสองบิตเป็น 1 \\
        ตัวอย่าง: $X$ \& $Y$ = $1011_2$ \& $0110_2$ = $0010_2$ \quad (ตรวจสอบทีละหลัก)

  \item \textbf{OR ($|$):} ผลลัพธ์เป็น 1 ก็ต่อเมื่อมีอย่างน้อยหนึ่งบิตเป็น 1 \\
        ตัวอย่าง: $X$ $|$ $Y$ = $1011_2$ $|$ $0110_2$ = $1111_2$

  \item \textbf{NOT ($\sim$):} กลับค่าทุกบิต (0 เป็น 1, 1 เป็น 0) \\
        ตัวอย่าง: $\sim X$ = $\sim 1011_2$ = $0100_2$

  \item \textbf{XOR (\textasciicircum):} ผลลัพธ์เป็น 1 ก็ต่อเมื่อทั้งสองบิตแตกต่างกัน \\
        ตัวอย่าง: $X$ \textasciicircum{} $Y$ = $1011_2$ \textasciicircum{} $0110_2$ = $1101_2$
\end{itemize}

\medskip
\textbf{ลำดับความสำคัญของตัวดำเนินการ (Precedence)} จากสูงไปต่ำ:

\begin{center}
\begin{tabular}{|c|c|c|}
\hline
\textbf{ลำดับ} & \textbf{ตัวดำเนินการ} & \textbf{ความหมาย} \\
\hline
1 (สูงสุด) & $\sim$ & NOT \\
2 & \& & AND \\
3 & \textasciicircum & XOR \\
4 (ต่ำสุด) & $|$ & OR \\
\hline
\end{tabular}
\end{center}

ตัวดำเนินการที่มีลำดับสูงกว่าจะถูกคำนวณก่อน วงเล็บ \textbf{( )} ใช้สำหรับจัดกลุ่มและจะถูกคำนวณก่อนตัวดำเนินการทั้งหมด

\medskip
\textbf{คำถาม:} กำหนดเลขฐานสองขนาด 4 หลักดังนี้

\begin{center}
$A = 1011_2$ \qquad $B = 0110_2$ \qquad $C = 1101_2$
\end{center}

จงหาผลลัพธ์ของนิพจน์ต่อไปนี้ \newline
(ตอบเป็นเลขฐานสอง \textbf{4 หลัก} โดยต้องเขียนเลข 0 นำหน้าด้วย เช่น $0011$ ไม่ใช่ $11$)

\begin{center}
\textbf{NOT((A AND B) XOR C) OR (NOT A AND B)}
\end{center}

\par\medskip
ตอบ ในกระดาษคำตอบ
% answer: 0100
% solution:
%   Step-by-step (all in 4-bit):
%   1) A AND B  = 1011 & 0110 = 0010
%   2) (A AND B) XOR C = 0010 ^ 1101 = 1111
%   3) NOT((A AND B) XOR C) = NOT(1111) = 0000
%   4) NOT A = NOT(1011) = 0100
%   5) NOT A AND B = 0100 & 0110 = 0100
%   6) 0000 OR 0100 = 0100
%   Answer: 0100₂ = 4₁₀
